\section{Introduction}
\label{sec:introduction}

Euclid's argument shows that there are infinitely many primes; it says
nothing about how densely they lie. Euler's observation~\cite{euler1737}
answers a sharper question: the primes are dense enough that the sum of
their reciprocals diverges, although the sum of the reciprocals of the
squares converges. In modern terms the result reads as follows; we state it
in \cref{thm:divergence} and prove it in \cref{sec:divergence}.

% The draft compared this with the twin primes; Brun's theorem belongs in
% a later version.
The main tool is the product formula of \cref{thm:euler-product}, which
writes the zeta function as a product over the primes. Its proof is a
careful version of the sieve, and it needs one fact about the integers,
recorded as \cref{assum:ufd}.\footnote{Unique factorisation is proved in
  every introductory text, e.g.\ \citet[Theorem~2]{hardywright}; see also
  \citep[Chapter~1, {\S}1.2]{primer2001} for a version that keeps track of
  the units.}

\paragraph{What is new here.} Nothing in the mathematics. The exposition
separates the two limits Euler took together, so that each step is an
inequality between finite sums; about 5\% of the text is spent on that
bookkeeping.

\begin{quote}
The sum of the reciprocals of the primes is infinite, and it grows like
the logarithm of the logarithm of its bound.
\end{quote}
This paraphrase of the result is the one usually quoted~\cite[p.~351]{primer2001}.
As the textbook puts it, ``the primes are not so sparse that their
reciprocals can be summed, though they come close''~\cite{primer2001}.

\vspace{1em}

The paper is organised as follows.
\begin{itemize}
  \item \Cref{sec:product} proves the product formula.
  \item \Cref{sec:divergence} deduces the divergence, in two steps:
  \begin{itemize}
    \item a lower bound for $\log \zeta(s)$ as $s \to 1^+$;
    \item an upper bound for the same quantity by the prime sum.
  \end{itemize}
  \item \Cref{app:proofs} holds the proofs we postponed, and
  \cref{sec:remarks,sec:introduction} frame them (page~\pageref{sec:remarks}).
\end{itemize}
A machine-readable copy of the argument is at
\url{https://example.org/euler%20primes}.

\newpage
